\begin{helper}
Let \(G\) and \(G'\) be finite \(\Delta\)-regular graphs, let \(\gamma\) be a
real parameter, and let \(\mu,\mu'\) be sampling laws satisfying
\(\noderef{RandomIndependentSetSampling}\) for \(G,\Delta,\gamma\) and
\(G',\Delta,\gamma\), respectively.  Let \(\nu\) and \(\nu'\) be
activation--priority laws on
\(\mathcal P(V)\times\mathbb R^V\) and
\(\mathcal P(V')\times\mathbb R^{V'}\) with the Bernoulli activation atom
formula and the finite priority lower-rectangle product formula at activation
ratio \(\gamma/\Delta\).

Let \(r\in V\), \(X\subseteq V\), and let \(\phi:V\to V'\) be injective.
Assume \(X\subseteq N_G(r)\) and that \(\phi\) preserves all adjacencies inside
\(X\cup\{r\}\).  For every outside incidence
\[
  x\in X,\qquad z\notin X\cup\{r\},\qquad xz\in E(G),
\]
let \(\beta(x,z)\in V'\) be a private replacement vertex.  Assume the
replacement vertices are distinct, fresh from \(\phi(X\cup\{r\})\), adjacent
to their intended vertex \(\phi(x)\), nonadjacent to \(\phi(r)\), adjacent to
no other \(\phi(y)\) with \(y\in X\), and that they exhaust all nonlocal
neighbors of each \(\phi(x)\).

Then there is one common actual hybrid witness system.  It consists of a
finite injective list \(z_0,\ldots,z_{n-1}\) of all outside blockers adjacent
to \(X\), candidate sets
\[
  B_k=\{x\in X:xz_k\in E(G)\},
\]
source common and extension coordinates, target common and extension
coordinates, a finite atom index set \(\Omega\), source and target cells
\(C_\omega,C'_\omega\), common nonnegative weights \(w(\omega)\), complete
source and target stage events \(S_j(\omega),T_j(\omega)\), and real
normalized stage masses \(P(j,\omega)\).  The coordinate clauses, cell-mass
clauses, pairwise-disjointness clauses, unit-support cover and null-complement
clauses, and the displayed formulas for \(S_j\) and \(T_j\) are exactly the
actual split-blocker clauses for this single witness surface.  Moreover
\[
  P(j,\omega)\ge0,\qquad w(\omega)=0\Rightarrow P(j,\omega)=0,
\]
the two endpoint cell identities hold,
\[
\operatorname{ofReal}(w(\omega)P(0,\omega))
 =\nu(\{(A,\pi):I_G(A,\pi)\cap X\ne\emptyset\}\cap C_\omega),
\]
\[
\operatorname{ofReal}(w(\omega)P(n,\omega))
 =\nu'(\{(A',\pi'):I_{G'}(A',\pi')\cap\phi(X)\ne\emptyset\}\cap C'_\omega),
\]
and every adjacent replacement step is monotone on every atom:
\[
  P(k,\omega)\le P(k+1,\omega)\qquad(k<n).
\]
The endpoint identities also aggregate over this same finite atom family:
\[
  \operatorname{ofReal}\!\left(\sum_{\omega\in\Omega}w(\omega)P(0,\omega)\right)
  =
  \nu\{(A,\pi):I_G(A,\pi)\cap X\ne\emptyset\},
\]
\[
  \operatorname{ofReal}\!\left(\sum_{\omega\in\Omega}w(\omega)P(n,\omega)\right)
  =
  \nu'\{(A',\pi'):I_{G'}(A',\pi')\cap\phi(X)\ne\emptyset\}.
\]
This node is the paper-facing bridge from the original split-blocker graph
interface and the two sampling laws to the aggregate witness system, including
the range reduction \(0\le\gamma/\Delta\le1\).  The lower
\(\noderef{SamplingSplitOutsideBlockerActualHybridWitnessAssembly}\) node has
the different role of constructing the actual coordinate/cell surface,
building the same-surface paired input on that surface, and assembling the
complete hybrid mass chain.
\end{helper}
\begin{proof}
First record the activation-ratio range needed by the actual assembly.  The
source sampling hypothesis is a \(\noderef{RandomIndependentSetSampling}\) law
and therefore gives \(\gamma>0\).  If \(\Delta=0\), then the Lean real
convention gives \(\gamma/(\Delta:\mathbb R)=0\), so the ratio lies in
\([0,1]\).  If \(0<\Delta\), the vertex \(r\) witnesses that \(V\) is
nonempty, and
\(\noderef{RandomIndependentSetSamplingGammaLeDelta}\) applied to the source
sampling law gives \(\gamma\le\Delta\).  Dividing by the positive real
\((\Delta:\mathbb R)\) gives
\[
  0\le \gamma/(\Delta:\mathbb R)\le 1 .
\]

Now apply
\(\noderef{SamplingSplitOutsideBlockerActualHybridWitnessAssembly}\) to the
graphs, the two activation--priority laws, the product atom and rectangle
formulas, the local-copy data \(\phi,r,X\), the private-blocker map \(\beta\),
and the split-blocker hypotheses.  That node first chooses the actual
coordinate package and prescribed common refined-cell package, then defines
the complete source and target stages and the adjacent boundary events on
that same data, constructs the paired mixed surface with its endpoint and
adjacent premises, and applies the complete-stage mass law.  Its conclusion
is exactly the witness system stated here: one list \(z_k\), one family
\(B_k\), one source/target coordinate
surface, one atom type \(\Omega\), one paired cell family
\(C_\omega,C'_\omega\), one common weight function \(w\), complete stages
\(S_j,T_j\), and one normalized mass family \(P(j,\omega)\) satisfying
nonnegativity, zero-weight normalization, the two endpoint cell identities,
and adjacent monotonicity.  No separate displayed-cell package, same-index
source--target equality, or outside-\(B_k\) residual atom equality is used.
The assembly node also supplies the endpoint aggregation over this same atom
family:
\[
  \operatorname{ofReal}\!\left(\sum_{\omega\in\Omega}w(\omega)P(0,\omega)\right)
  =
  \nu\{(A,\pi):I_G(A,\pi)\cap X\ne\emptyset\},
\]
\[
  \operatorname{ofReal}\!\left(\sum_{\omega\in\Omega}w(\omega)P(n,\omega)\right)
  =
  \nu'\{(A',\pi'):I_{G'}(A',\pi')\cap\phi(X)\ne\emptyset\}.
\]
\end{proof}
